% The complete Feller median/symmetrization route, with atom-aware formulas.
\section{Feller's weak-law presentation: proofs}
\label{sec:annex-p15-feller-weak-law-presentation}

Use Chapter~\ref{sec:p15-feller-weak-law-presentation} for statements
and conventions. The source passages are Feller V.2, V.5, V.(6.1)
and VII.7, Theorem 1
\cite[pp.~137, 148--150, 234--236]{Feller1971ProbabilityII}.

\begin{annexitem}{annex:p15.feller-symmetrization-structure}{thm:p15-feller-symmetrization-structure}{Medians, product copies and the distribution formulas}
\begin{proof}
The CDF limits at the two infinities show that
\(\{x:F(x)\geq1/2\}\) is nonempty and bounded below. Let \(m\)
be its infimum. For every \(h>0\), monotonicity and the defining
infimum give \(F(m+h)\geq1/2\). Right continuity yields
\(F(m)\geq1/2\). For every \(x<m\), \(F(x)<1/2\);
taking the left limit gives \(F(m-)\leq1/2\).
Thus \(m\) is a finite median. The identities
\(\mathbb P(X-m\leq0)=F(m)\) and
\(\mathbb P(X-m\geq0)=1-F(m-)\) prove the centering assertion.
The median need not be unique: every point of \([-1,1]\) is a
median for the fair law on \(\{-1,1\}\).

For reflection,
\[
 \mathbb P(-X\leq x)=\mathbb P(X\geq-x)
                     =1-\mathbb P(X<-x)=1-F((-x)-).
\]
At a continuity point of \(F\) at \(-x\) this reduces to the
source's \(1-F(-x)\).
Construct copies on \(\Omega^2\) with the finite product
probability \(\mathbb Q=\mathbb P\otimes\mathbb P\).
Coordinate maps are measurable and their laws are the original laws.
Every finite family of constraints on copied observations is a
rectangle; its probability factors between the two coordinates,
and original independence factors the constraints within each
coordinate. Thus both copied sequences together are independent
whenever the original sequence is independent. Differences from
disjoint copied pairs remain independent, and common input laws
give common difference laws.

For one pair, the product law and Tonelli applied to the Borel event
\(\{(u,v):u-v\leq x\}\) give
\[
 {}^{\circ}F(x)
 =\int_{\mathbb R}\int_{\mathbb R}\mathbf1_{\{u\leq x+v\}}
                       F\{du\}F\{dv\}
 =\int_{\mathbb R}F(x+v)\,F\{dv\}.
\]
This is also the convolution law of \(X+(-X')\), as in
Theorem~\ref{thm:p15-convolution-law}.
Interchanging the two coordinates preserves product probability:
it preserves rectangle probabilities, hence the full product measure
by uniqueness. This interchange negates the difference, proving
its symmetry. Applying the reflection identity to this symmetric
law gives its exact left-limit CDF relation, including jumps.
Finally \((X-b)-(X'-b)=X-X'\) pointwise for each deterministic
finite \(b\), so symmetrization is translation invariant.
The same identity applies separately to each deterministically
centered partial sum.
\end{proof}
\end{annexitem}

\begin{annexitem}{annex:p15.feller-symmetrization-inequalities}{thm:p15-feller-symmetrization-inequalities}{Every source tail inequality and the median specialization}
\begin{proof}
If both \(|X|\leq t/2\) and \(|X'|\leq t/2\), their
difference has magnitude at most \(t\). The union bound and equal
laws therefore give
\[
 \mathbb P(|X-X'|>t)
 \leq\mathbb P(|X|>t/2)+\mathbb P(|X'|>t/2)
 =2\mathbb P(|X|>t/2).
\]
For the lower bound, the two disjoint events
\[
 \{X>t+a,\ X'\leq a\},\qquad
 \{X<-t-a,\ X'\geq-a\}
\]
each force \(|X-X'|>t\). Independence gives their respective
probabilities as products. The assumed lower bounds for the two
copy probabilities consequently give
\[
 \mathbb P(|X-X'|>t)
 \geq p\{\mathbb P(X>t+a)+\mathbb P(X<-t-a)\}
 =p\,\mathbb P(|X|>t+a).
\]
Here \(a\geq0\) and \(t>0\), so the two tail events are disjoint.
The cases \(p=0\) and \(p=1\) cause no division or exception.
If zero is a median, take \(a=0,p=1/2\) to obtain V.(5.8).

For any median \(m\), \(|m|\geq m\) and \(-|m|\leq m\).
Thus
\[
 \mathbb P(X\leq|m|)\geq\mathbb P(X\leq m)\geq1/2,\qquad
 \mathbb P(X\geq-|m|)\geq\mathbb P(X\geq m)\geq1/2.
\]
Taking the nonnegative radius \(|m|\) and \(p=1/2\) in V.(5.7)
gives the stated median tail comparison. Finally, apply V.(5.6)
to \(X-b,X'-b\) at threshold \(2t\) and use translation
invariance of their difference. This gives the translated upper
bound used for centered partial sums.
\end{proof}
\end{annexitem}

\begin{annexitem}{annex:p15.feller-largest-term-inequality}{thm:p15-feller-largest-term-inequality}{The first-largest-term proof and both threshold conventions}
\begin{proof}
The vector \((Y_1,\ldots,Y_n)\) has a product law on
\(\mathbb R^n\). Reflection of any chosen coordinates preserves
this law, because each coordinate law is symmetric. This follows
first on rectangles and then on their generated sigma-algebra.
For each \(j\), the event \(\{J=j\}\) is specified by
\[
 |Y_i|<|Y_j|\ (i<j),\qquad |Y_i|\leq|Y_j|\ (i>j).
\]
It is Borel measurable and is unchanged by any coordinate reflection.
Consequently \(J,M,R\) are measurable.

On \(\{J=j\}\), reflecting coordinate \(j\) changes \((M,R)\)
to \((-M,R)\) and preserves both the product law and the event
\(\{J=j\}\). Reflecting all coordinates other than \(j\)
changes it to \((M,-R)\) with the same preservation properties.
For any Borel set of pairs, sum the resulting equal probabilities
over the finite disjoint partition \(\{J=j\}\).
This proves invariance of the joint law under either sign change,
and hence under both. It proves symmetry of \(M\) and of \(S=M+R\)
without claiming that \(M\) and \(R\) are independent.

For \(t>0\), joint symmetry under changing the sign of \(R\)
gives
\[
 \begin{aligned}
 \mathbb P(M>t,R\geq0)
 &=\tfrac12\mathbb P(M>t,R\ne0)+\mathbb P(M>t,R=0)\\
 &\geq\tfrac12\mathbb P(M>t).
 \end{aligned}
\]
The zero-remainder mass is included on the nonnegative side.
On this event \(S=M+R>t\), so
\(\mathbb P(S>t)\geq\frac12\mathbb P(M>t)\).
Since \(S\) and \(M\) are symmetric and \(t>0\), their positive
and negative tail probabilities are equal and disjoint. Multiplying
the one-sided comparison by two gives
\[
 \mathbb P(|S|>t)\geq\tfrac12\mathbb P(|M|>t)
                    =\tfrac12\mathbb P(\max_i|Y_i|>t).
\]
The same reasoning with \(M\geq t\) gives the non-strict
comparison. Positivity of \(t\) keeps the two absolute-value tails
disjoint in that case too. This is the source's largest-term proof.

For common coordinate law \(G\), independence gives the exact
maximum probabilities
\[
 \begin{aligned}
 \mathbb P(\max_i|Y_i|\leq t)
   &=[G(t)-G((-t)-)]^n=(1-q_>(t))^n,\\
 \mathbb P(\max_i|Y_i|<t)
   &=[G(t-)-G(-t)]^n=(1-q_\geq(t))^n.
 \end{aligned}
\]
At continuity points these reduce to the power appearing in source
V.(5.13). For \(0\leq q\leq1\), \(1-q\leq e^{-q}\),
so the two maximal comparisons yield the stated exponential
bounds, including \(q=0\) and \(q=1\).
Finally
\[
 1-G(t)+G(-t)
   =\mathbb P(Y_1>t)+\mathbb P(Y_1\leq-t)
   \leq\mathbb P(|Y_1|\geq t).
\]
Monotonicity of \(q\mapsto1-e^{-nq}\) proves the source CDF form
at every threshold, including atoms. Its left side is non-strict;
the separate strict form uses the strict absolute-value tail.
\end{proof}
\end{annexitem}

\begin{annexitem}{annex:p15.feller-integration-by-parts}{thm:p15-feller-integration-by-parts}{The complete partition and telescoping argument}
\begin{proof}
Set \(M=\sup_{[a,b]}|u'|<\infty\).
Partition \(a=x_0<\cdots<x_N=b\) into equal intervals of length
\(h=(b-a)/N\), and put
\(\Delta F_j=F(x_j)-F(x_{j-1})\),
\(\Delta u_j=u(x_j)-u(x_{j-1})\).
Define
\[
 A=\int_{(a,b]}[u(b)-u(x)]\,F\{dx\},\qquad
 B=\int_a^b u'(x)[F(x)-F(a)]\,dx.
\]
The finite approximations
\[
 A_N=\sum_{j=1}^N[u(b)-u(x_j)]\Delta F_j,\qquad
 B_N=\sum_{j=1}^N\Delta u_j[F(x_{j-1})-F(a)]
\]
are equal. Indeed,
\[
 \begin{aligned}
 A_N
 &=\sum_{j=1}^N\Delta F_j\sum_{\ell=j+1}^N\Delta u_\ell\\
 &=\sum_{\ell=1}^N\Delta u_\ell
                         \sum_{j=1}^{\ell-1}\Delta F_j
 =B_N.
 \end{aligned}
\]
For \(x\in(x_{j-1},x_j]\), the mean-value bound gives
\(|u(x)-u(x_j)|\leq Mh\). Therefore
\[
 |A-A_N|\leq Mh\sum_j\Delta F_j\leq Mh.
\]
On the Lebesgue interval \([x_{j-1},x_j]\),
\(0\leq F(x)-F(x_{j-1})\leq\Delta F_j\).
The fundamental theorem of calculus identifies
\(\Delta u_j=\int_{x_{j-1}}^{x_j}u'(x)\,dx\), so
\[
 |B-B_N|
 \leq\sum_j\int_{x_{j-1}}^{x_j}
           |u'(x)|\,|F(x)-F(x_{j-1})|\,dx
 \leq Mh\sum_j\Delta F_j\leq Mh.
\]
Hence \(|A-B|\leq2Mh\to0\), proving \(A=B\).
Expanding these two finite integrals gives
\[
 u(b)[F(b)-F(a)]-\int_{(a,b]}u(x)\,F\{dx\}
  =\int_a^b u'(x)F(x)\,dx-F(a)[u(b)-u(a)].
\]
Rearranging proves the stated Stieltjes formula. The partition uses
\((x_{j-1},x_j]\), so the right endpoint atoms and exclusion of
the left endpoint are preserved throughout.
\end{proof}
\end{annexitem}

\begin{annexitem}{annex:p15.feller-tail-moment-equivalence}{thm:p15-feller-tail-moment-equivalence}{The moment identity, endpoint atoms and its full converse}
\begin{proof}
Apply integration by parts to \(u(x)=x^2\) on \([-t,t]\):
\[
 \int_{(-t,t]}x^2\,F\{dx\}
 =t^2[F(t)-F(-t)]-2\int_{-t}^t xF(x)\,dx.
\]
Reflecting the negative half of the Lebesgue integral gives
\[
 \int_{-t}^t xF(x)\,dx
   =\int_0^t y[F(y)-F(-y)]\,dy
   =\int_0^t[y-\tau(y)]\,dy.
\]
Since \(t^2[F(t)-F(-t)]=t^2-t\tau(t)\), substitution and
division by \(t>0\) give
\(\sigma_h(t)=-\tau(t)+(2/t)\int_0^t\tau(y)\,dy\).
The only point added in passing to the closed interval is \(-t\).
Its second-moment contribution divided by \(t\) is
\(t\mathbb P(X=-t)\), and its first-moment contribution is
\(-t\mathbb P(X=-t)\). This proves both other exact identities.

Write \(W=|X|\) and \(r(t)=t\mathbb P(W>t)\). The disjoint
tail definitions imply, for \(t>0\),
\[
 r(t)\leq\tau(t)\leq t\mathbb P(W\geq t)
                  \leq t\mathbb P(W>t/2)=2r(t/2).
\]
Thus the first two vanishing conditions are equivalent.
Also \(0\leq\tau(t)\leq t\), and
\(t\mathbb P(X=-t)\leq\tau(t)\). The moment identities give
\[
 0\leq\sigma(t)\leq\frac2t\int_0^t\tau(y)\,dy.
\]
If \(\tau(t)\to0\), choose \(A\) such that \(\tau(y)\leq\varepsilon\)
for \(y\geq A\). For \(t>A\), the last average is at most
\((2/t)\int_0^A\tau(y)\,dy+2\varepsilon\).
Its finite head tends to zero, and \(\varepsilon\) is arbitrary.
This proves \(\sigma(t)\to0\).

For the converse, partition \(\{W>t\}\) into the disjoint events
\(\{2^kt<W\leq2^{k+1}t\}\), \(k\geq0\).
The variable \(W\) is finite, so they cover that event.
For each \(k\),
\[
 \begin{aligned}
 t\,\mathbb P(2^kt<W\leq2^{k+1}t)
 &\leq\frac{t}{(2^kt)^2}
                  \mathbb E[W^2\mathbf1_{\{W\leq2^{k+1}t\}}]\\
 &=2^{1-k}\sigma(2^{k+1}t).
 \end{aligned}
\]
Summing gives
\[
 r(t)\leq4\sup_{s\geq2t}\sigma(s).
\]
If \(\sigma(s)\to0\), the right side tends to zero, so \(r(t)\to0\)
and hence \(\tau(t)\to0\).
Finally \(\sigma_h(t)\leq\sigma(t)\), while
\([-t,t]\subset(-2t,2t]\) implies
\(\sigma(t)\leq2\sigma_h(2t)\).
These comparisons prove equivalence of the two moment conditions.
All four conditions are therefore equivalent, with no omitted
converse referral.
\end{proof}
\end{annexitem}

\begin{annexitem}{annex:p15.feller-truncation-inequality}{thm:p15-feller-truncation-inequality}{The source two-error estimate and its normalization}
\begin{proof}
The event \(\{|S_n-m_n'|>n\varepsilon\}\) is contained in
\[
 \{S_n\ne S_n'\}\ \cup
                         \{|S_n'-m_n'|>n\varepsilon\}.
\]
The first event is contained in the union of the events
\(\{|X_k|>s\}\), whose total probability is at most
\(n\mathbb P(|X_1|>s)\).
The truncated variables are independent, identically distributed
and bounded by \(s\). Their sum has mean \(m_n'=nm(s)\) and
variance \(n\operatorname{Var}(X_1')\).
Chebyshev therefore bounds the second probability by
\[
 \frac{n\operatorname{Var}(X_1')}{n^2\varepsilon^2}
       \leq\frac{\mathbb E[(X_1')^2]}{n\varepsilon^2}.
\]
Adding the two probabilities proves the claim.
Only the truncated variables needed expectations or variances.
The scale \(n^{-1}\) follows from retaining all \(n\) variance
contributions in the numerator before division by \(n^2\).
\end{proof}
\end{annexitem}

\begin{annexitem}{annex:p15.feller-generalized-weak-law}{thm:p15-feller-generalized-weak-law}{The full median chain and the source sufficiency proof}
\begin{proof}
Assume first that \(\tau(t)\to0\). Choose the closed truncated
centers \(\mu_n=m(n)\), and use the truncation inequality with
\(s=n\). For every \(\varepsilon>0\),
\[
 \begin{aligned}
 \mathbb P(|S_n/n-\mu_n|>\varepsilon)
 &\leq\varepsilon^{-2}\sigma(n)+n\mathbb P(|X_1|>n)\\
 &\leq\varepsilon^{-2}\sigma(n)+\tau(n)\longrightarrow0.
 \end{aligned}
\]
The moment equivalence established \(\sigma(n)\to0\), so every
term in the source sufficiency argument is justified.
The half-open center differs by
\(m_h(n)-m(n)=n\mathbb P(X_1=-n)\), which is at most \(\tau(n)\).
It therefore gives the same concentration conclusion.

Conversely, assume real deterministic \(\mu_n\) satisfy the stated
concentration property. Construct two independent copies of the entire
original sequence on the finite product probability space, with law
\(\mathbb Q=\mathbb P\otimes\mathbb P\).
Denote the first-coordinate copies again by \(X_i\), and write
\(X_i'\) for the second-coordinate copies, \(S_n'=\sum_iX_i'\),
\({}^{\circ}X_i=X_i-X_i'\), and
\({}^{\circ}S_n=S_n-S_n'=\sum_i{}^{\circ}X_i\).
The structural theorem proves that these differences are i.i.d.
and symmetric. They are the direct symmetrizations required by
Feller's argument, not an arbitrary symmetric replacement.
The variables \(S_n-n\mu_n\) and \(S_n'-n\mu_n\) are independent
with identical laws; their difference is \({}^{\circ}S_n\).
Apply V.(5.6) to this translated pair:
\[
 2\mathbb P(|S_n-n\mu_n|>n\varepsilon)
       \geq\mathbb Q(|{}^{\circ}S_n|>2n\varepsilon).
\]

Let \(a\) be any finite median of the common input law, whose
existence was proved above. The strict form of the largest-term
exponential inequality gives
\[
 \mathbb Q(|{}^{\circ}S_n|>2n\varepsilon)
 \geq\tfrac12\left(1-
        e^{-n\mathbb Q(|{}^{\circ}X_1|>2n\varepsilon)}\right).
\]
The median specialization of V.(5.7), with radius \(|a|\) and
probability bound \(1/2\), gives
\[
 \mathbb Q(|{}^{\circ}X_1|>2n\varepsilon)
       \geq\tfrac12\mathbb P(|X_1|>2n\varepsilon+|a|).
\]
The exponential expression is increasing in this tail probability.
Combining the three inequalities proves the exact displayed
median chain in the theorem.

Its left side tends to zero by assumption. Hence, for every fixed
\(\varepsilon>0\),
\[
 n\mathbb P(|X_1|>2n\varepsilon+|a|)\longrightarrow0.
\]
Indeed, a subsequence on which this nonnegative quantity were at
least a positive constant would keep the final exponential bound
away from zero. Take \(\varepsilon=1/4\). For all sufficiently
large \(n\), \(n/2+|a|\leq n\), so
\(n\mathbb P(|X_1|>n)\to0\).
For real \(t\to\infty\), put \(n=\lfloor t\rfloor\); then
\[
 t\mathbb P(|X_1|>t)
       \leq(n+1)\mathbb P(|X_1|>n)\longrightarrow0.
\]
The proven equivalence between this tail limit and \(\tau(t)\to0\)
completes necessity. The argument uses the source's median and
largest-term comparisons throughout; no first-crossing maximal
inequality or central limit theorem is substituted.
\end{proof}
\end{annexitem}

\begin{annexitem}{annex:p15.feller-centering-corollaries}{thm:p15-feller-centering-corollaries}{Permissible centers and the two source specializations}
\begin{proof}
The theorem makes the canonical center \(m(n)\) work whenever
\(\tau(t)\to0\). If \(c_n-m(n)\to0\), a triangle-inequality
split at \(\varepsilon/2\) makes \(c_n\) work too.
Conversely, if both centers work, then for \(\varepsilon>0\)
the event that their deterministic distance exceeds \(\varepsilon\)
forces at least one of their random errors to exceed
\(\varepsilon/2\). Thus
\[
 \mathbf1_{\{|c_n-m(n)|>\varepsilon\}}
 \leq\mathbb P(|S_n/n-c_n|>\varepsilon/2)
       +\mathbb P(|S_n/n-m(n)|>\varepsilon/2)\longrightarrow0.
\]
The deterministic indicator is eventually zero, proving convergence
of the center difference.

If \(W=|X_1|\) is integrable, then
\[
 0\leq\tau(t)\leq t\mathbb P(W\geq t)
                    \leq\mathbb E[W\mathbf1_{\{W\geq t\}}]\longrightarrow0
\]
by dominated convergence. The same theorem gives
\(m(n)=\mathbb E[X_1\mathbf1_{\{W\leq n\}}]\to\mathbb EX_1\).
The canonical weak law and the vanishing deterministic center
difference give the finite-mean conclusion.
If \(F\) is symmetric, the bounded random variable
\(X_1\mathbf1_{\{|X_1|\leq n\}}\) has the same law as its
negative. Its expectation is therefore its own negative and equals
zero. This proves the symmetric-law conclusion without asserting
existence of the untruncated expectation.
\end{proof}
\end{annexitem}

\begin{annexitem}{annex:p15.feller-symmetrization-examples}{thm:p15-feller-symmetrization-examples}{The density computations and the distinction between symmetry and symmetrization}
\begin{proof}
If \(X\) has density \(f\), the independent difference has density
\(h(z)=\int_{\mathbb R}f(z+y)f(y)\,dy\).
This is the already proved convolution-density formula, using the
reflected density for \(-X'\) and a one-variable affine substitution.
For \(f(y)=\lambda e^{-\lambda y}\mathbf1_{\{y\geq0\}}\),
\(\lambda>0\), it gives
\[
 h(z)=\lambda^2e^{-\lambda z}
       \int_{\max(0,-z)}^\infty e^{-2\lambda y}\,dy
       =\frac{\lambda}{2}e^{-\lambda|z|}.
\]
This nonnegative density integrates to one by splitting at zero.
For the uniform density on \((0,1)\), the same integral is the
length of \((0,1)\cap(-z,1-z)\), namely \(1-|z|\) for
\(|z|<1\), and zero otherwise. Its integral over the line is one;
the displayed value at the two endpoints is also zero.

For Bernoulli \(X\), list the four independent pairs \((X,X')\).
The differences \(1,-1\) each have probability \(p(1-p)\),
and zero has probability \(p^2+(1-p)^2\). These add to one.
In particular
\({}^{\circ}F(0)=1-p(1-p)\) and
\({}^{\circ}F(0-)=p(1-p)\), so their sum is one, as required
by the exact symmetry relation. This includes \(p=0,1\).

For any real probability law and any \(\delta>0\), partition
\(\mathbb R\) into countably many half-open intervals of length
\(\delta/2\). At least one interval has positive probability;
otherwise countable additivity would give total mass zero.
Independent copies both lie in that interval with positive
probability, and their distance is then less than \(\delta\).
Thus every directly symmetrized law has positive probability
inside \((-\delta,\delta)\). The fair law on \(\{-1,1\}\)
is symmetric but gives that interval zero probability when
\(\delta=1\), proving that it cannot be a copied-difference law.

For \(X=-1\), the CDF gives
\(1-F(t)+F(-t)=1\) for \(0<t\leq1\) and zero for \(t>1\).
The closed truncated second moment includes the atom exactly when
\(t\geq1\), whereas the half-open interval includes it exactly
when \(t>1\). Dividing those moments by \(t\) proves the three
stated formulas. At \(t=1\), the closed moment is one and the
half-open moment is zero; their difference is exactly
\(t\mathbb P(X=-t)=1\).
\end{proof}
\end{annexitem}

\begin{annexitem}{annex:p15.feller-proof-boundaries}{thm:p15-feller-proof-boundaries}{Constants, atoms and the explicit failure models}
\begin{proof}
For Bernoulli probability \(0<p<1\) and \(0<t<1\),
\(\mathbb P(|X-X'|>t)=2p(1-p)\) and
\(\mathbb P(|X|>t/2)=p\). Their ratio tends to two as
\(p\downarrow0\), so no smaller universal upper constant works.
For independent fair signs \(X,X'\), zero is a median,
\(\mathbb P(|X|>t)=1\) for \(0<t<1\), and
\(\mathbb P(|X-X'|>t)=1/2\).
This attains the median lower constant.
For independent fair signs \(Y_1,Y_2\), their maximum absolute
value is one, while their sum is zero with probability \(1/2\)
and has magnitude two otherwise. This attains the largest-term
constant at the same thresholds.

If \(Y_1\) is a fair sign and \(Y_2=0\), then \(M=Y_1\)
and \(R=0\). For \(0<t<1\),
\(\mathbb P(M>t,R\geq0)=1/2\) and
\(\mathbb P(M>t,R<0)=0\), while the at-least-half comparison
used in the proof remains true.
For a fair Bernoulli variable used as both \(X\) and \(X'\),
zero is a median but the difference is zero; for \(0<t<1\)
the proposed median lower side would be
\((1/2)\mathbb P(|X|>t)=1/4\). This verifies failure without
independence. For \(Y_2=-Y_1\) with a fair sign, the sum is
zero and the maximum magnitude is one, disproving the largest-term
claim for dependent symmetric variables. Independent constants
\(1,-1\) give the same failure after individual symmetry is removed.

The constant law \(X=1\) has independent constant copies and zero
copied difference, but \(\mathbb P(|X|>t)=1\) when \(0<t<1\).
Zero is not a median, so the unqualified zero-median inequality
would indeed fail. For a fair symmetric sign with CDF \(G\),
at \(t=1\) the source CDF tail is
\(1-G(1)+G(-1)=1/2\).
Its exponential lower bound is positive, while
\(\mathbb P(|Y_1|>1)=0\). The actual non-strict left side is
one. Thus the two threshold versions cannot be interchanged at
that atom.

For two independent fair signs, \(S_2/2\) is \(-1,0,1\)
with probabilities \(1/4,1/2,1/4\). With cutoff \(s=1\)
the truncation changes nothing and has second moment one.
At \(\varepsilon=9/10\) the exact deviation probability is
\(1/2\). The corrected single-summand variance bound is
\(1/(2(9/10)^2)=50/81\), whereas \(1/(4(9/10)^2)=25/81\)
is smaller than the actual probability.

Finally, for the two-sided logarithmic dyadic law, the identity
\[
 \sum_{j\geq1}q^j/j
       =\int_0^q(1-t)^{-1}\,dt=-\log(1-q)
                         \quad(0\leq q<1)
\]
follows by integrating the nonnegative geometric series.
At \(q=1/2\) it proves normalization of the displayed masses.
Each part expectation is
\((2\log2)^{-1}\sum_{j\geq1}1/j=\infty\).
For \(2^k\leq t<2^{k+1}\),
\[
 \mathbb P(|X|>t)
  =\frac1{\log2}\sum_{j\geq k+1}\frac1{j2^j}
  \leq\frac{2^{-k}}{(k+1)\log2}.
\]
Thus \(t\mathbb P(|X|>t)\to0\), and the proved tail equivalence
gives \(\tau(t)\to0\). Symmetry makes every closed truncated
mean zero. Feller's theorem as proved here therefore gives the
zero weak limit, while the untruncated expectation is undefined.
This example uses the newly reconstructed source route, without
an appeal to the earlier independent weak-law proof.
\end{proof}
\end{annexitem}
